\documentclass[24pt, a3paper, portrait]{tikzposter}%Set paper, basic font size and orientation
\usepackage[utf8]{inputenc}
\usepackage{blindtext}
\usepackage{fontspec}
\usepackage[super]{natbib}%My preferred bibliography package, another can be used
\usepackage[english]{babel}
\usepackage{jabbrv}%For abbreviating journal names

%------------------------Style Definitions-----------------------

\setmainfont{Arial}%NHS font, use Frutiger if you have the license and the font files: https://www.overleaf.com/blog/using-your-favourite-fonts-with-sharelatex-2013-04-02
\newcommand{\tSup}[1]{\textsuperscript{#1}}%A shortcut for superscripts

\definecolorstyle{christieColourStyle} {%Colour styling. Defines colours and where to use them
    %https://www.england.nhs.uk/nhsidentity/identity-guidelines/colours/
    \definecolor{NHSBlue}{RGB}{0,94,184}%NHS blue
    \definecolor{NHSDarkBlue}{RGB}{0,48,135}%NHS dark blue
    \definecolor{NHSDarkPink}{RGB}{124,40,85}%NHS dark pink
}{
    %BG
    \colorlet{backgroundcolor}{white}
    \colorlet{framecolor}{white}
    %Title
    \colorlet{titlefgcolor}{NHSBlue}
    \colorlet{titlebgcolor}{white}
    \colorlet{authorfgcolor}{NHSDarkPink}
    %Blocks
    \colorlet{blocktitlebgcolor}{white}
    \colorlet{blocktitlefgcolor}{NHSBlue}
    \colorlet{blockbodybgcolor}{white}
    \colorlet{blockbodyfgcolor}{black}
    %Innerblocks
    \colorlet{innerblocktitlebgcolor}{NHSDarkPink}
    \colorlet{innerblocktitlefgcolor}{white}
    \colorlet{innerblockbodybgcolor}{NHSBlue}
    \colorlet{innerblockbodyfgcolor}{white}
    %Notes
    \colorlet{notefgcolor}{NHSDarkPink}
    \colorlet{notebgcolor}{white}
    \colorlet{noteframecolor}{NHSDarkPink}
}

\settitle {%Define the title section layout
    \@titlegraphic \par
    \vspace*{1em}
    \color{titlefgcolor} {\bfseries \Huge \@title \par}
    \vspace*{0.75em}
    \color{authorfgcolor} {\Large \@author \par}
    \vspace*{0.75em}
    \color{titlefgcolor} {\large\@institute}
}

%--------------------------Title Content-------------------------

\title{Christie Poster Template}
\author{Frank Brewster\tSup{1}, Barry White\tSup{2}}
\date{May 2021}
\institute{\tSup{1}The Christie NHS Foundation Trust, Manchester, UK; \tSup{2}The University of Manchester, Manchester, UK}
%A second uni logo can be added to the top using this:
%\titlegraphic{\includegraphics[height=5cm]{Uni logo file name} \hfill \includegraphics[height=5cm]{christieLogoTight.jpg}}
\titlegraphic{\hfill \includegraphics[height=5cm]{christieLogoTight.jpg}}

\usecolorstyle{christieColourStyle}

%----------------------------------------------------------------
%--------------------------END PREAMBLE--------------------------
%----------------------------------------------------------------

\begin{document}

\maketitle[width=0.95\linewidth]

\block{Abstract}
{
    \blindtext
}

\begin{columns}%Number of columns is flexible. Different parts of the page can have different numbers of columns by starting a new columns environment
    \column{0.33}%The size of this column
    \block{Introduction}{\blindtext}%A section
    
    \block{Methodology}{\blindtext 
        \innerblock{An inner-block}{Some inner block text to provide some emphasis on a key point within another block.}
    \blindtext}
    
    \column{0.33}
    
    \block{Results}{I like figure \ref{fig:dummy}, it is a very nice figure.\\ \blindtext}
    
    \block{}{%A figure
        \begin{tikzfigure}[some figure caption\label{fig:dummy}]
            \includegraphics[width=\linewidth]{Large_The_Christie_embrace_RGB.jpg}
        \end{tikzfigure}
    }
    
    \block{Discussion}{A study demonstrated this\cite{VanHerk2003}.}
    
    \column{0.33}
    \block{Conclusion}{\blindtext}
    \block{}{
        \bibliographystyle{fbPosterShort}%A custom bib style with 1 author and lots of abbreviations. Journal abbreviations by jabbrv.
        \bibliography{bib.bib}%Your bib file
    }
    
\end{columns}

%Bottom embrace logo
\node[above right,outer sep = 0.025\linewidth] at (bottomleft) {\includegraphics[height=5cm]{Large_The_Christie_embrace_RGB.jpg}};

\end{document}
